\documentclass[a4paper]{arrowhead}

\usepackage[yyyymmdd]{datetime}
\usepackage{etoolbox}
\usepackage[utf8]{inputenc}
\usepackage[T1]{fontenc}
\usepackage{multirow}
\usepackage{longtable}
\usepackage{tikz}
\usepackage{adjustbox}
\AtBeginEnvironment{longtable}{\small}
\BeforeBeginEnvironment{tikzpicture}{\begin{adjustbox}{max width=\textwidth}}
\AfterEndEnvironment{tikzpicture}{\end{adjustbox}}
\usetikzlibrary{positioning,shapes.multipart,arrows.meta}

\renewcommand{\dateseparator}{-}
\setlength{\emergencystretch}{3em}
\lstset{language=HTTP,breakatwhitespace=false,columns=fullflexible}

\newcommand{\fparam}[1]{\textit{\textcolor{ArrowheadBlue}{#1}}}

%% Special references
\newcommand{\mref}[2][]{{\textcolor{ArrowheadPurple}{\hyperref[sec:model:#2]{\ifstrempty{#1}{#2}{#1}}}}}
\newcommand{\oref}[1]{{\textcolor{ArrowheadBlue}{\hyperref[sec:operations:#1]{#1}}}}
\newcommand{\pdef}[1]{{\textcolor{ArrowheadGrey}{#1 \label{sec:model:primitives:#1} \label{sec:model:primitives:#1s}}}}
\newcommand{\pref}[1]{{\textcolor{ArrowheadGrey}{\hyperref[sec:model:primitives:#1]{#1}}}}

% as template_idd.tex; the input and output arguments are passed already formatted (\mref or \pref)
\newrobustcmd\osubsection[5]{
  \addtocounter{subsection}{1}
  \addcontentsline{toc}{subsection}{\protect\numberline{\thesubsection}operation \textcolor{ArrowheadBlue}{#1}}
  \subsection*{
    \thesubsection\quad
    #2 \textcolor{ArrowheadPurple}{#3} \\
    \small
    \hspace*{0.075\textwidth}\begin{minipage}{0.1\textwidth}
      \vspace*{1mm}
      Operation: \\
      \notblank{#4}{Input: \\}{}
      \notblank{#5}{Output: \\}{}
    \end{minipage}%
    \begin{minipage}{0.825\textwidth}
      \vspace*{1mm}
      \textcolor{ArrowheadBlue}{#1} \\
      \notblank{#4}{#4 \\}{}
      \notblank{#5}{#5 \\}{}
    \end{minipage}
  }
  \label{sec:operations:#1}
}
\newrobustcmd\msubsection[2]{
  \addtocounter{subsection}{1}
  \addcontentsline{toc}{subsection}{\protect\numberline{\thesubsection}#1 \textcolor{ArrowheadPurple}{#2}}
  \subsection*{\thesubsection\quad#1 \textcolor{ArrowheadPurple}{#2}}
  \label{sec:model:#2}
}
%%

\begin{document}
\sloppy

%% Arrowhead Document Properties
\ArrowheadTitle{serviceOrchestration MQTT/TLS/JSON}
\ArrowheadServiceID{serviceOrchestration}
\ArrowheadType{Interface Design Description}
\ArrowheadTypeShort{IDD}
\ArrowheadVersion{5.2.2}
\ArrowheadDate{2026-10-22}
\ArrowheadAuthor{Jerker Delsing}
\ArrowheadStatus{RELEASE CANDIDATE}
\ArrowheadContact{jerker.delsing@ltu.se}
\ArrowheadFooter{\href{www.arrowhead.eu}{www.arrowhead.eu}}
\ArrowheadSetup
%%

%% Front Page
\begin{center}
  \vspace*{1cm}
  \huge{\arrowtitle}

  \vspace*{0.2cm}
  \LARGE{\arrowtype}
  \vspace*{1cm}

  \Large{Service ID: \textit{"\arrowid"}}

  \vspace*{\fill}

  % Front Page Image

  \vspace*{1cm}
  \vspace*{\fill}

  % Front Page Abstract
  \begin{abstract}
    This document describes the MQTT protocol with TLS payload security and JSON payload encoding variant of the \texttt{service\allowbreak{}Orchestration} service of the SimpleStoreServiceOrchestration system.
  \end{abstract}

  \vspace*{1cm}
\end{center}

\newpage
%%

%% Table of Contents
\tableofcontents
\newpage
%%

\section{Overview}
\label{sec:overview}

This document describes the \texttt{service\allowbreak{}Orchestration} service interface of the SimpleStoreServiceOrchestration system over MQTT, which provides service orchestration: a consumer asks which provider it should use for a service, and receives the matching, authorized providers with the interfaces to call.

It's implemented using protocol, encoding as stated in below table:
\begin{table}[ht!]
  \centering
  \begin{tabular}{|l|l|l|}
    \hline
    \rowcolor{gray!33} Profile type & Type & Version \\ \hline
    Transfer protocol & MQTT & 3.1.1 \cite{mqtt311} \\ \hline
    Data encryption & TLS & 1.3 \cite{rfc8446}, 1.2 \cite{rfc5246} \\ \hline
    Encoding & JSON & RFC 8259 \cite{rfc8259} \\ \hline
    Compression & N/A & - \\ \hline
    Semantics & N/A & - \\ \hline
    Ontology & N/A & - \\ \hline
  \end{tabular}
  \caption{Communication and semantics details used for the \texttt{service\allowbreak{}Orchestration} service interface}
  \label{tab:comunication_semantics_profile}
\end{table}

This document provides the Interface Design Description IDD to the \textit{serviceOrchestration -- Service Description} document \cite{ssso-serviceOrchestration-sd}.
The operations, their semantics and the data models are the same as in the HTTP/TLS/JSON IDD \cite{ssso-serviceOrchestration-http-idd}; this document specifies the MQTT binding.

The MQTT interface is offered only when the system runs with \texttt{mqtt.api.enabled=true}; it is then registered with the interface template \texttt{generic\_\allowbreak{}mqtts} (\texttt{generic\_\allowbreak{}mqtt} without TLS), the broker address and port (\texttt{mqtt.\allowbreak{}broker.\allowbreak{}address}, \texttt{mqtt.\allowbreak{}broker.\allowbreak{}port}) and the base topic \texttt{arrowhead/\allowbreak{}serviceorchestration/\allowbreak{}orchestration/\allowbreak{}}. In the Arrowhead-test-cloud-java deployment the MQTT interfaces of the core systems are disabled.

\newpage

\section{Service Operations}
\label{sec:operations}

The provider subscribes to one topic per operation, the base topic followed by the operation name:

\noindent\begin{longtable}{| p{5cm} | p{10.4cm} |} \hline
\rowcolor{gray!33} Operation & Request topic \\ \hline \endhead
\texttt{pull} & \texttt{arrowhead/\allowbreak{}serviceorchestration/\allowbreak{}orchestration/\allowbreak{}pull} \\ \hline
\texttt{subscribe} & \texttt{arrowhead/\allowbreak{}serviceorchestration/\allowbreak{}orchestration/\allowbreak{}subscribe} \\ \hline
\texttt{unsubscribe} & \texttt{arrowhead/\allowbreak{}serviceorchestration/\allowbreak{}orchestration/\allowbreak{}unsubscribe} \\ \hline
\end{longtable}

A consumer publishes a request envelope (\mref{MqttRequestTemplate}) to the topic and receives the answer (\mref{MqttResponseTemplate}) on its own \texttt{responseTopic}. The \texttt{authentication} field carries the requester's identity according to the provider's authentication policy: \texttt{IDENTITY-TOKEN//<token>} (outsourced or internal policy), \texttt{SYSTEM//<name>} (declared policy) or the requester's X.509 certificate (certificate policy). The \texttt{status} of the response uses the HTTP status code numbers (200, 201, 204, 400, 401, 403, 404, 408, 423, 500, 503); on error the payload is an \mref{ErrorMessage}.

\begin{lstlisting}[caption={Request envelope.}]
{
  "traceId": "a1b2c3",
  "authentication": "IDENTITY-TOKEN//<token>",
  "responseTopic": "arrowhead/response/TemperatureDisplay/a1b2c3",
  "qosRequirement": 1,
  "params": {},
  "payload": {}
}
\end{lstlisting}

\begin{lstlisting}[caption={Response envelope.}]
{
  "status": 200,
  "traceId": "a1b2c3",
  "receiver": "TemperatureDisplay",
  "payload": {}
}
\end{lstlisting}

\osubsection{pull}{PUBLISH}{\texttt{arrowhead/\allowbreak{}serviceorchestration/\allowbreak{}orchestration/\allowbreak{}pull}}{\mref{OrchestrationRequest}}{\mref{OrchestrationResponse}}

Returns the store entries of the requester for the requested service definition in priority order; preferred providers first (\texttt{ONLY\_\allowbreak{}PREFERRED}: only them); with \texttt{MATCHMAKING} only the best entry per service definition.

\begin{lstlisting}[caption={Payload of an \oref{pull} request.}]
{
  "serviceRequirement": {
    "serviceDefinition": "temperatureReading",
    "operations": [
      "FILTER"
    ],
    "versions": [
      "1.0.0"
    ],
    "alivesAt": "2026-10-22T12:00:00Z",
    "metadataRequirements": [
      {
        "unit": "celsius"
      }
    ],
    "interfaceTemplateNames": [
      "string"
    ],
    "interfaceAddressTypes": [
      "string"
    ],
    "interfacePropertyRequirements": [
      {
        "unit": "celsius"
      }
    ],
    "securityPolicies": [
      "string"
    ],
    "preferredProviders": [
      "string"
    ]
  },
  "orchestrationFlags": {
    "MATCHMAKING": true
  },
  "qualityRequirements": [
    {
      "type": "PING",
      "operation": "FILTER",
      "requirements": {
        "maxLatency": 100
      }
    }
  ],
  "exclusivityDuration": 60
}
\end{lstlisting}

\osubsection{subscribe}{PUBLISH}{\texttt{arrowhead/\allowbreak{}serviceorchestration/\allowbreak{}orchestration/\allowbreak{}subscribe}}{\mref[Orchestration\allowbreak{}Subscription\allowbreak{}Request]{OrchestrationSubscriptionRequest}}{\pref{String}}

Creates (or replaces) a push subscription of the requester for itself; with \texttt{trigger=\allowbreak{}true} the orchestration is executed and pushed immediately. Returns the subscription identifier, 201 for a new subscription.

Optional entries of \texttt{params}: \texttt{trigger}.

\begin{lstlisting}[caption={Payload of an \oref{subscribe} request.}]
{
  "targetSystemName": "TemperatureDisplay",
  "orchestrationRequest": {
    "serviceRequirement": {
      "serviceDefinition": "temperatureReading",
      "operations": [
        "FILTER"
      ],
      "versions": [
        "1.0.0"
      ],
      "alivesAt": "2026-10-22T12:00:00Z",
      "metadataRequirements": [
        {
          "unit": "celsius"
        }
      ],
      "interfaceTemplateNames": [
        "string"
      ],
      "interfaceAddressTypes": [
        "string"
      ],
      "interfacePropertyRequirements": [
        {
          "unit": "celsius"
        }
      ],
      "securityPolicies": [
        "string"
      ],
      "preferredProviders": [
        "string"
      ]
    },
    "orchestrationFlags": {
      "MATCHMAKING": true
    },
    "qualityRequirements": [
      {
        "type": "PING",
        "operation": "FILTER",
        "requirements": {
          "maxLatency": 100
        }
      }
    ],
    "exclusivityDuration": 60
  },
  "notifyInterface": {
    "protocol": "HTTPS",
    "properties": {
      "address": "192.168.0.20",
      "port": "8080",
      "method": "POST",
      "path": "/orchestration/notify"
    }
  },
  "duration": 3600
}
\end{lstlisting}

\osubsection{unsubscribe}{PUBLISH}{\texttt{arrowhead/\allowbreak{}serviceorchestration/\allowbreak{}orchestration/\allowbreak{}unsubscribe}}{\pref{String}}{}

Removes a subscription of the requester. Returns 200 if removed, 204 if it did not exist.

\begin{lstlisting}[caption={Payload of an \oref{unsubscribe} request.}]
"3fa85f64-5717-4562-b3fc-2c963f66afa6"
\end{lstlisting}

\newpage

\section{Data Models}
\label{sec:model}

Here, all data objects that can be part of the service calls associated with this service are listed in alphabetic order.
Each \textit{struct} denotes a JSON \pref{Object} with the named fields; fields that are not needed may be omitted.

\msubsection{struct}{MqttRequestTemplate}

The request envelope.

\noindent\begin{tabularx}{\textwidth}{| p{4.25cm} | p{3.5cm} | X |} \hline
\rowcolor{gray!33} Object Field & Value Type & Description \\ \hline
"traceId" & \pref{String} & Identifier chosen by the requester, returned in the response. \\ \hline
"authentication" & \pref{String} & Identity of the requester (see Section \ref{sec:operations}). \\ \hline
"responseTopic" & \pref{String} & Topic on which the response is published. \\ \hline
"qosRequirement" & \pref{Number} & MQTT QoS of the response (0, 1 or 2). \\ \hline
"params" & \pref{Object}$<$\pref{String}$>$ & Optional parameters (the query parameters of the HTTP binding). \\ \hline
"payload" & \pref{Object} & The operation input. \\ \hline
\end{tabularx}

\msubsection{struct}{MqttResponseTemplate}

The response envelope.

\noindent\begin{tabularx}{\textwidth}{| p{4.25cm} | p{3.5cm} | X |} \hline
\rowcolor{gray!33} Object Field & Value Type & Description \\ \hline
"status" & \pref{Number} & Status code. \\ \hline
"traceId" & \pref{String} & The trace identifier of the request. \\ \hline
"receiver" & \pref{String} & Name of the requester. \\ \hline
"payload" & \pref{Object} & The operation output, or an \mref{ErrorMessage}. \\ \hline
\end{tabularx}

\msubsection{struct}{ErrorMessage}




\noindent\begin{tabularx}{\textwidth}{| >{\raggedright\arraybackslash}p{3.9cm} | >{\raggedright\arraybackslash\small}p{5.3cm} | X |} \hline
\rowcolor{gray!33} Object Field & Value Type & Description \\ \hline
"error\allowbreak{}Message" & \pref{String} &  \\ \hline
"error\allowbreak{}Code" & \pref{Number} &  \\ \hline
"exception\allowbreak{}Type" & \pref{ExceptionType} &  \\ \hline
"origin" & \pref{String} &  \\ \hline
\end{tabularx}



\msubsection{struct}{AuthorizationTokenGenerationResponse}

A generated token.


\noindent\begin{tabularx}{\textwidth}{| >{\raggedright\arraybackslash}p{3.9cm} | >{\raggedright\arraybackslash\small}p{5.3cm} | X |} \hline
\rowcolor{gray!33} Object Field & Value Type & Description \\ \hline
"token\allowbreak{}Type" & \pref{AuthorizationTokenType} & Token type: \texttt{TIME\_\allowbreak{}LIMITED\_\allowbreak{}TOKEN}, \texttt{USAGE\_\allowbreak{}LIMITED\_\allowbreak{}TOKEN}, \texttt{SELF\_\allowbreak{}CONTAINED\_\allowbreak{}TOKEN} or \texttt{TRANSLATION\_\allowbreak{}BRIDGE\_\allowbreak{}TOKEN}. \\ \hline
"target\allowbreak{}Type" & \pref{AuthorizationTargetType} & Target type: \texttt{SERVICE\_\allowbreak{}DEF} or \texttt{EVENT\_\allowbreak{}TYPE}. \\ \hline
"token" & \pref{String} & The token (opaque string, Base64 or JSON Web Token). \\ \hline
"usage\allowbreak{}Limit" & \pref{Number} & Number of allowed uses (usage-limited tokens). \\ \hline
"expires\allowbreak{}At" & \pref{DateTime} & UTC expiry of time-limited tokens. \\ \hline
\end{tabularx}



\msubsection{struct}{MetadataRequirement}

A JSON object of metadata requirements: each key is a metadata path (dot notation), each value the required value or an operator object such as \texttt{\{"op": "LESS\_\allowbreak{}THAN",\allowbreak{} "value": 10\}}.



\msubsection{struct}{OrchestrationNotifyInterface}

Notification endpoint of push orchestration.


\noindent\begin{tabularx}{\textwidth}{| >{\raggedright\arraybackslash}p{3.9cm} | >{\raggedright\arraybackslash\small}p{5.3cm} | X |} \hline
\rowcolor{gray!33} Object Field & Value Type & Description \\ \hline
"protocol" & \pref{String} & Protocol of the notification: \texttt{HTTP}, \texttt{HTTPS}, \texttt{MQTT} or \texttt{MQTTS}. \\ \hline
"properties" & \pref{Object}$<$\allowbreak{}\pref{String}$>$ & For HTTP(S) \texttt{address}, \texttt{port}, \texttt{method}, \texttt{path}; for MQTT(S) \texttt{topic} (and broker \texttt{address}, \texttt{port}). \\ \hline
\end{tabularx}



\msubsection{struct}{OrchestrationRequest}

An orchestration request.


\noindent\begin{tabularx}{\textwidth}{| >{\raggedright\arraybackslash}p{3.9cm} | >{\raggedright\arraybackslash\small}p{5.3cm} | X |} \hline
\rowcolor{gray!33} Object Field & Value Type & Description \\ \hline
"service\allowbreak{}Requirement" & \mref[Orchestration\allowbreak{}Service\allowbreak{}Requirement]{OrchestrationServiceRequirement} & Requirements on the service the consumer wants to use. \\ \hline
"orchestration\allowbreak{}Flags" & \pref{Object}$<$\allowbreak{}\pref{Boolean}$>$ & Orchestration flags, e.g. \texttt{MATCHMAKING}, \texttt{ONLY\_\allowbreak{}PREFERRED}, \texttt{ALLOW\_\allowbreak{}INTERCLOUD}. \\ \hline
"quality\allowbreak{}Requirements" & \pref{Array}$<$\allowbreak{}\mref{QoSRequirement}$>$ & Quality-of-service requirements (evaluated by Device\allowbreak{}Qo\allowbreak{}SEvaluator). \\ \hline
"exclusivity\allowbreak{}Duration" & \pref{Number} & Requested exclusive reservation of the provider in seconds. \\ \hline
\end{tabularx}



\msubsection{struct}{OrchestrationResponse}

The orchestration result.


\noindent\begin{tabularx}{\textwidth}{| >{\raggedright\arraybackslash}p{3.9cm} | >{\raggedright\arraybackslash\small}p{5.3cm} | X |} \hline
\rowcolor{gray!33} Object Field & Value Type & Description \\ \hline
"results" & \pref{Array}$<$\allowbreak{}\mref{OrchestrationResult}$>$ & The orchestration results, best first. \\ \hline
"warnings" & \pref{Array}$<$\allowbreak{}\pref{String}$>$ & Warnings of the orchestration, e.g. \texttt{autoMatchmaking}, \texttt{notExclusive}, \texttt{interCloud}. \\ \hline
\end{tabularx}



\msubsection{struct}{OrchestrationResult}

One provider of the requested service.


\noindent\begin{tabularx}{\textwidth}{| >{\raggedright\arraybackslash}p{3.9cm} | >{\raggedright\arraybackslash\small}p{5.3cm} | X |} \hline
\rowcolor{gray!33} Object Field & Value Type & Description \\ \hline
"service\allowbreak{}Instance\allowbreak{}Id" & \pref{String} & Service instance identifier \texttt{\textless{}System\textgreater{}\textbar{}\allowbreak{}\textless{}service\allowbreak{}Definition\textgreater{}\textbar{}\allowbreak{}\textless{}version\textgreater{}}. \\ \hline
"cloud\allowbreak{}Identitifer" & \pref{String} & Cloud identifier of the provider (\texttt{LOCAL} for the own cloud). The JSON name keeps the upstream spelling. \\ \hline
"provider\allowbreak{}Name" & \pref{String} & Name of the provider system. \\ \hline
"service\allowbreak{}Definitition" & \pref{String} & Service definition name. The JSON name keeps the upstream spelling. \\ \hline
"version" & \pref{Version} & Version of the service instance. \\ \hline
"alive\allowbreak{}Until" & \pref{String} & Time until which the service instance is registered (its expiry), if any. \\ \hline
"exclusive\allowbreak{}Until" & \pref{String} & Time until which the instance is reserved exclusively for the consumer, if reserved. \\ \hline
"metadata" & \pref{Object} & Metadata of the service instance. \\ \hline
"interfaces" & \pref{Array}$<$\allowbreak{}\mref[Service\allowbreak{}Instance\allowbreak{}Interface\allowbreak{}Response]{ServiceInstanceInterfaceResponse}$>$ & Interfaces the consumer can use (only matching ones). \\ \hline
"authorization\allowbreak{}Tokens" & \pref{Object}$<$\allowbreak{}\pref{Object}$<$\allowbreak{}\mref[Authorization\allowbreak{}Token\allowbreak{}Generation\allowbreak{}Response]{AuthorizationTokenGenerationResponse}$>$$>$ & Authorization tokens generated for the consumer, per interface template and operation. \\ \hline
\end{tabularx}



\msubsection{struct}{OrchestrationServiceRequirement}

Requirements on the requested service.


\noindent\begin{tabularx}{\textwidth}{| >{\raggedright\arraybackslash}p{3.9cm} | >{\raggedright\arraybackslash\small}p{5.3cm} | X |} \hline
\rowcolor{gray!33} Object Field & Value Type & Description \\ \hline
"service\allowbreak{}Definition" & \pref{String} & Service definition name. \\ \hline
"operations" & \pref{Array}$<$\allowbreak{}\pref{String}$>$ & Service operations the consumer wants to use. \\ \hline
"versions" & \pref{Array}$<$\allowbreak{}\pref{Version}$>$ & Accepted service versions. \\ \hline
"alives\allowbreak{}At" & \pref{DateTime} & Only instances that are still registered (not expired) at this UTC time are matched. \\ \hline
"metadata\allowbreak{}Requirements" & \pref{Array}$<$\allowbreak{}\mref{MetadataRequirement}$>$ & Requirements on the service instance metadata; each map is one alternative. \\ \hline
"interface\allowbreak{}Template\allowbreak{}Names" & \pref{Array}$<$\allowbreak{}\pref{String}$>$ & Accepted interface templates, e.g. \texttt{generic\_\allowbreak{}https}. \\ \hline
"interface\allowbreak{}Address\allowbreak{}Types" & \pref{Array}$<$\allowbreak{}\pref{String}$>$ & Address types the provider interface must offer. \\ \hline
"interface\allowbreak{}Property\allowbreak{}Requirements" & \pref{Array}$<$\allowbreak{}\mref{MetadataRequirement}$>$ & Requirements on the interface properties; each map is one alternative. \\ \hline
"security\allowbreak{}Policies" & \pref{Array}$<$\allowbreak{}\pref{String}$>$ & Accepted interface security policies. \\ \hline
"preferred\allowbreak{}Providers" & \pref{Array}$<$\allowbreak{}\pref{String}$>$ & Provider systems to prefer. \\ \hline
\end{tabularx}



\msubsection{struct}{OrchestrationSubscriptionRequest}

A push orchestration subscription.


\noindent\begin{tabularx}{\textwidth}{| >{\raggedright\arraybackslash}p{3.9cm} | >{\raggedright\arraybackslash\small}p{5.3cm} | X |} \hline
\rowcolor{gray!33} Object Field & Value Type & Description \\ \hline
"target\allowbreak{}System\allowbreak{}Name" & \pref{String} & Consumer system; default the requester. \\ \hline
"orchestration\allowbreak{}Request" & \mref{OrchestrationRequest} & The orchestration request executed for the subscription. \\ \hline
"notify\allowbreak{}Interface" & \mref[Orchestration\allowbreak{}Notify\allowbreak{}Interface]{OrchestrationNotifyInterface} & Where and how the orchestration result is pushed to the target system. \\ \hline
"duration" & \pref{Number} & Requested lifetime of the subscription in seconds; no value means unlimited. \\ \hline
\end{tabularx}



\msubsection{struct}{QoSRequirement}

One quality-of-service requirement.


\noindent\begin{tabularx}{\textwidth}{| >{\raggedright\arraybackslash}p{3.9cm} | >{\raggedright\arraybackslash\small}p{5.3cm} | X |} \hline
\rowcolor{gray!33} Object Field & Value Type & Description \\ \hline
"type" & \pref{String} & Type of the QoS evaluation, as known by Device\allowbreak{}Qo\allowbreak{}SEvaluator. \\ \hline
"operation" & \pref{String} & QoS operation: \texttt{FILTER} or \texttt{SORT}. \\ \hline
"requirements" & \pref{Object} & Parameters of the QoS requirement. \\ \hline
\end{tabularx}



\msubsection{struct}{ServiceInstanceInterfaceResponse}

One interface of a service instance.


\noindent\begin{tabularx}{\textwidth}{| >{\raggedright\arraybackslash}p{3.9cm} | >{\raggedright\arraybackslash\small}p{5.3cm} | X |} \hline
\rowcolor{gray!33} Object Field & Value Type & Description \\ \hline
"template\allowbreak{}Name" & \pref{String} & Name of the interface template, snake\_case (e.g. \texttt{generic\_\allowbreak{}https}). \\ \hline
"protocol" & \pref{String} & Protocol of the interface, e.g. \texttt{http}, \texttt{https}, \texttt{mqtt}, \texttt{mqtts}. \\ \hline
"policy" & \pref{String} & Security policy of the interface: \texttt{NONE}, \texttt{CERT\_\allowbreak{}AUTH} or one of the token policies. \\ \hline
"properties" & \pref{Object} & Interface properties defined by the template, e.g. \texttt{accessAddresses}, \texttt{accessPort}, \texttt{basePath}, \texttt{operations}. \\ \hline
\end{tabularx}



\subsection{Primitives}
\label{sec:model:primitives}

As all messages are encoded using the JSON format \cite{rfc8259}, the following primitive constructs, part of that standard, become available.

\noindent\begin{tabularx}{\textwidth}{| p{4.6cm} | X |} \hline
\rowcolor{gray!33} Type & Description \\ \hline
\pdef{Array} & A JSON Array \cite{rfc8259}; \pref{Array}$<$A$>$ has elements of type A. \\ \hline
\pdef{AuthorizationTargetType} & One of \texttt{SERVICE\_\allowbreak{}DEF}, \texttt{EVENT\_\allowbreak{}TYPE} (JSON String). \\ \hline
\pdef{AuthorizationTokenType} & One of \texttt{TIME\_\allowbreak{}LIMITED\_\allowbreak{}TOKEN}, \texttt{USAGE\_\allowbreak{}LIMITED\_\allowbreak{}TOKEN}, \texttt{SELF\_\allowbreak{}CONTAINED\_\allowbreak{}TOKEN}, \texttt{TRANSLATION\_\allowbreak{}BRIDGE\_\allowbreak{}TOKEN} (JSON String). \\ \hline
\pdef{Boolean} & JSON Boolean, \texttt{true} or \texttt{false}. \\ \hline
\pdef{DateTime} & JSON String with an ISO 8601 UTC timestamp, e.g. \texttt{2026-10-22T12:00:00Z}. \\ \hline
\pdef{ExceptionType} & One of \texttt{ARROWHEAD}, \texttt{INVALID\_\allowbreak{}PARAMETER}, \texttt{AUTH}, \texttt{FORBIDDEN}, \texttt{DATA\_\allowbreak{}NOT\_\allowbreak{}FOUND}, \texttt{TIMEOUT}, \texttt{LOCKED}, \texttt{INTERNAL\_\allowbreak{}SERVER\_\allowbreak{}ERROR}, \texttt{EXTERNAL\_\allowbreak{}SERVER\_\allowbreak{}ERROR} (JSON String). \\ \hline
\pdef{Number} & JSON Number, an integer. \\ \hline
\pdef{Object} & JSON Object; \pref{Object}$<$A$>$ has string keys and values of type A. \\ \hline
\pdef{String} & JSON String. \\ \hline
\pdef{Version} & JSON String, semantic version \texttt{major.minor.patch}. \\ \hline
\end{tabularx}

\newpage

\bibliographystyle{IEEEtran}
\bibliography{bibliography}

\newpage

\section{Revision History}
\subsection{Amendments}

\noindent\begin{tabularx}{\textwidth}{| p{1cm} | p{3cm} | p{2cm} | X | p{4cm} |} \hline
\rowcolor{gray!33} No. & Date & Version & Subject of Amendments & Author \\ \hline

1 & 2026-10-22 & \arrowversion & Document created from the Arrowhead LaTeX templates for the 5.2.2 Java Spring implementation, including the local extensions of this workspace. & Jerker Delsing \\ \hline
\end{tabularx}

\subsection{Quality Assurance}

\noindent\begin{tabularx}{\textwidth}{| p{1cm} | p{3cm} | p{2cm} | X |} \hline
\rowcolor{gray!33} No. & Date & Version & Approved by \\ \hline

1 & & \arrowversion  &  \\ \hline

\end{tabularx}

\end{document}
